\documentclass[12pt]{article}
\usepackage{../../ficha}

\begin{document}

\ficha{Números Complexos na Forma Algébrica}[Operações, Conjugado e Equações]

\instrucoes{Apresenta os resultados na forma algébrica $a + bi$, com $a, b \in \R$. Lembra-te: $i^2 = -1$.}

\secao{Parte Real, Parte Imaginária e Potências de $i$}{Resolve.}
\begin{exercicios}
  \item Indica a parte real e a parte imaginária de $z = 3 - 5i$.
  \item Indica a parte real e a parte imaginária de $z = -2i$ e de $w = 7$.
  \item Calcula $i^3$, $i^4$, $i^{10}$ e $i^{2027}$.
  \item Calcula $i + i^2 + i^3 + i^4$.
  \item Determina $k \in \R$ de modo que $z = (k - 2) + (k + 3)i$ seja um imaginário puro.
  \item Determina $k \in \R$ de modo que $z = (k^2 - 1) + (k - 2)i$ seja um número real não nulo.
  \item Determina $a, b \in \R$ tais que $(a + 1) + (2b - 3)i = 4 + 5i$.
\end{exercicios}

\secao{Adição, Subtração e Multiplicação}{Calcula.}
\begin{exercicios}
  \item $(3 + 2i) + (1 - 5i)$
  \item $(4 - i) - (-2 + 3i)$
  \item $2(1 + 3i) - 3(2 - i)$
  \item $(2 + i)(3 - 2i)$
  \item $(1 + i)^2$
  \item $(1 - i)^2$
  \item $(3 + 4i)(3 - 4i)$
  \item $i(2 - 5i)$
  \item $(2 - i)^3$
  \item $(1 + i)^{8}$ \; (usa o resultado do exercício 5)
\end{exercicios}

\secao{Conjugado, Módulo e Divisão}{Calcula.}
\begin{exercicios}
  \item O conjugado de $z = 5 - 2i$ e o valor de $z \cdot \overline{z}$.
  \item $|3 - 4i|$
  \item $|{-5i}|$
  \item $\dfrac{1}{i}$
  \item $\dfrac{2 + i}{1 - i}$
  \item $\dfrac{3 - 2i}{2 + 3i}$
  \item $\dfrac{5}{1 + 2i}$
  \item $\dfrac{(1 + i)^2}{1 - i}$
  \item Mostra que $\overline{z + w} = \overline{z} + \overline{w}$ para $z = 1 + 2i$ e $w = 3 - i$.
  \item Mostra que $z + \overline{z}$ é sempre um número real e $z - \overline{z}$ é sempre um imaginário puro.
\end{exercicios}

\secao{Equações em $\C$}{Resolve em $\C$.}
\begin{exercicios}
  \item $z^2 + 9 = 0$
  \item $z^2 + 2z + 5 = 0$
  \item $z^2 - 4z + 13 = 0$
  \item $(1 + i)z = 3 - i$
  \item $2z - i = 4 + 3i$
  \item $z + 2\overline{z} = 6 - i$ \; (escreve $z = a + bi$)
  \item $z^2 = 2i$ \; (escreve $z = a + bi$)
  \item $z^3 - 1 = 0$ \; (fatoriza $z^3 - 1$)
\end{exercicios}

\secao{Desafios}{Resolve os problemas.}
\begin{exercicios}
  \item Determina um polinómio do 2.º grau com coeficientes reais cujas raízes sejam $2 + 3i$ e $2 - 3i$.
  \item Mostra que, se um polinómio de coeficientes reais tem a raiz $z$, então também tem a raiz $\overline{z}$.
  \item Sabendo que $1 + i$ é raiz de $z^3 - 4z^2 + 6z - 4$, determina as outras raízes.
  \item Determina os números complexos $z$ tais que $z^2 = \overline{z}$.
\end{exercicios}

\end{document}
